\documentclass[11pt,a4paper]{article}
\usepackage[UTF8]{ctex}
\usepackage{geometry}
\geometry{left=2.5cm,right=2.5cm,top=2.5cm,bottom=2.5cm}
\usepackage{amsmath,amssymb,amsthm}
\usepackage{hyperref}
\usepackage{booktabs}
\usepackage{enumitem}
\usepackage{xltabular}
\usepackage{array}
\hypersetup{colorlinks=true,linkcolor=blue,urlcolor=blue}

\title{安全厂商告警降噪技术方案调查报告}
\author{基于公开信息的系统性梳理}
\date{\today}

\begin{document}

\maketitle

\begin{center}
\textit{
本报告系统梳理了当前国内主流安全厂商在告警降噪领域的技术方案、产品形态与效果数据。\\
旨在为“基于生成论的网络安全分析能力调度系统”提供与工程实践对话的实证基础。\\
所有数据均来源于厂商公开披露信息及国家级测试结果。
}
\end{center}

\vspace{5mm}

\begin{abstract}
告警洪流是安全运营领域最普遍的痛点之一。近年来，国内主流安全厂商纷纷推出基于AI大模型的告警降噪方案，形成了各具特色的技术路线。本报告系统梳理了奇安信、深信服、腾讯云、阿里云、微步在线、绿盟科技、360、亚信安全等八家厂商的告警降噪技术方案，涵盖其技术架构、降噪效果、产品形态及公开测试表现。在此基础上，对各方案的共性技术路径与差异化策略进行横向对比分析，并指出其与“够用就行”理论体系的潜在对话接口。本报告旨在为理论体系与工程实践的对话搭建桥梁。
\end{abstract}

\tableofcontents

\section{引言：告警降噪的市场背景}

安全运营中心（SOC）面临的告警洪流问题已非个别现象。据奇安信对万人规模以上企业的调查，86\%的企业安全运营人员不到10人，研判比例不足5\%；13\%的企业运营人员有10-30人，研判比例在5\%-10\%；仅有1\%的企业配备了30人以上的安全运营团队，但研判告警比例也仅达到10\%。这些数据说明，人力不足导致大量告警被忽视，已成为网络安全当前面临的最大漏洞。

在此背景下，国内主流安全厂商纷纷将AI大模型技术引入告警降噪领域，形成了以“AI大模型+安全平台”为核心范式的技术路线。2024年和2025年，国家互联网应急中心（CNCERT）连续两年组织“人工智能技术赋能网络安全应用测试”，其中“网络安全告警日志降噪”均为核心测试场景。这表明告警降噪已成为AI赋能安全运营的标杆性应用场景。

\section{厂商技术方案详述}

\subsection{奇安信：智能分诊与“NGSOC+N”联动}

奇安信在告警降噪领域提出了“智能告警分诊”概念，由\textbf{专家经验模型和人工智能模型共同驱动}，对来自不同设备的告警信息实现自动过滤误报告警、自动识别精准告警、自动分辨重要告警、自动合并重复告警、自动处置无效告警，并重新分类和优先级排序。

\textbf{技术架构}：奇安信的告警降噪方案以NGSOC（态势感知与安全运营平台）为核心，深度融合QAX-GPT安全大模型，形成“AI+SOC智能安全运营方案”。其核心技术组件包括：
\begin{itemize}
    \item \textbf{大数据处理引擎（NOAH）}：负责海量告警数据的采集、存储与预处理；
    \item \textbf{大数据关联引擎（SABRE）}：负责多源告警的关联分析与事件聚合；
    \item \textbf{安全编排与自动化响应引擎（SOAR）}：负责自动化处置流程的执行；
    \item \textbf{AI大模型（QAX-GPT）}：负责智能研判与辅助决策。
\end{itemize}

在实际运营中，奇安信采用\textbf{“双重告警降噪”}机制：第一层通过NGSOC平台的“智能分诊”功能将告警筛选过程压缩到秒级，在海量告警中识别出约1\%的重要告警；第二层通过QAX-GPT安全机器人对初步过滤的告警进行“确诊”，进一步筛选出攻击成功的有效告警。这种“NGSOC+N”的联动机制（NGSOC+天眼+QAX-GPT+SOAR）可将原本需要数小时完成的告警筛选过程压缩到秒级。

\textbf{效果数据}：
\begin{itemize}
    \item 实验室环境下告警降噪率可达90\%左右；
    \item 预置分诊模型可消除98\%以上的告警噪声；
    \item 在实战攻防演练中，通过“智能分诊+GPT机器人”双重降噪，可消除90\%的告警噪声；
    \item NGSOC实现1分钟事件定性，5分钟完成影响面评估。
\end{itemize}

\textbf{产品形态}：奇安信天眼威胁监测与分析系统（NDR产品）已为全球6000余家客户提供服务。NGSOC平台预置2400+解析规则、1200+关联规则、100+分诊模型，支持1000+种数据源的自动化解析。

\subsection{深信服：XDR+安全GPT双引擎}

深信服的告警降噪方案以\textbf{“XDR+安全GPT”双引擎}为核心架构。XDR平台负责多源数据的融合分析与告警聚合，安全GPT负责智能研判与自动化值守。

\textbf{技术架构}：
\begin{itemize}
    \item \textbf{XStream技术}：实现多源数据的融合分析；
    \item \textbf{威胁定性能力}：对告警进行一键定性（业务误触、普通病毒、自动化攻击、定向攻击）；
    \item \textbf{情境检测分析}：高质量精准还原攻击故事线；
    \item \textbf{安全GPT}：化身虚拟安全专家，7×24小时自主值守。
\end{itemize}

\textbf{效果数据}：
\begin{itemize}
    \item 告警降噪率达95\%，事件研判效率提升10倍；
    \item 日均告警从原有的上千万条削减到10+条；
    \item 在某保险客户POC中，单日XDR平台接收各网端日志4000万+条，生成安全告警6.1万条，GPT降噪后单日仅个位数条需分析关注的事件，整体降噪率超过99\%；
    \item GPT分析研判准确率：内网事件93.1\%，互联网事件94.8\%；
    \item 在安全GPT辅助下，一个驻场安服12小时工作可完成至少2880条告警研判/天，效率提升近10倍；
    \item 2024年度国家攻防演练中平均告警降噪率达99\%，准确率达98\%，自动化处置率超80\%。
\end{itemize}

\textbf{产品形态}：深信服XDR已服务1000+用户，其AI安全运营方案覆盖“预防-检测-研判-响应”全周期智能闭环。

\subsection{腾讯云：Multi-Agent告警研判体系}

腾讯云鼎实验室采用\textbf{Multi-Agent（多智能体）分工协作架构}进行告警自动研判。

\textbf{技术架构——五层告警研判体系}：
\begin{enumerate}[label=第\arabic*层：]
    \item \textbf{告警数据源}：接入主机安全、容器安全、云防火墙、WAF等多源告警；
    \item \textbf{告警触发与归并}：筛选未阻断的危险告警进行触发，按关键维度（如同源IP+同目的IP+同用户名）归并；
    \item \textbf{Multi-Agent研判引擎}：采用Plan-React-Analysis架构，将复杂的告警分析拆成“规划→调查→研判”三步，多个Agent分工协作；
    \item \textbf{MCP工具层}：集成60+腾讯云安全工具，覆盖情报查询、资产画像、日志分析等；
    \item \textbf{输出层}：输出结论+证据链+处置建议。
\end{enumerate}

以“异常登录”告警为例，Agent会从7个维度展开调查：威胁情报、攻击历史、登录基线、登录后行为、告警关联、漏洞风险、资产画像。

\textbf{效果数据}：
\begin{itemize}
    \item 告警降噪率提升至97.7\%；
    \item 以某MMORPG游戏公司为例：7天原始告警17万条，经同特征聚合后1834条，Agent研判后真正需关注的仅42条；
    \item Agent能从一条“攻击失败”告警中追溯出2个月前已失陷的主机。
\end{itemize}

\subsection{阿里云：告警分级与事件降噪规则}

阿里云的告警降噪方案以\textbf{通知策略与事件管理}为核心，提供系统化的告警降噪能力。

\textbf{技术架构}：
\begin{itemize}
    \item \textbf{通知策略}：作为CloudMonitor 2.0告警中心的核心调度单元，负责将告警规则与外部集成产生的可观测事件进行订阅、降噪、路由、生命周期管理与自动化行动；
    \item \textbf{合并降噪}：通过分组字段将相似事件聚合在一起，减少通知疲劳度；
    \item \textbf{分级管理}：通过云监控、SLS日志服务和EventBridge等产品，实现告警的分级管理，确保关键威胁优先响应。
\end{itemize}

\textbf{效果数据}：据公开信息，阿里云告警降噪方案可减少无效告警量约60\%-80\%，通过分组、去重、抑制、静默等降噪规则实现。

\textbf{产品形态}：告警降噪能力内置于云监控、日志服务、事件中心等产品中。

\subsection{微步在线：三维IP画像与威胁情报降噪}

微步在线的告警降噪方案以\textbf{威胁情报为核心驱动力}，通过“TI+AI”技术内核实现告警降噪。

\textbf{技术架构}：
\begin{itemize}
    \item \textbf{三维IP画像技术}：从深度、广度和活跃度三个维度为告警信息赋予更多上下文；
    \item \textbf{威胁情报平台NGTIP}：提供威胁情报、漏洞情报和态势情报三重服务；
    \item \textbf{深度联动}：联动各类安全设备，实现威胁自动化闭环处置。
\end{itemize}

\textbf{效果数据}：
\begin{itemize}
    \item 三维IP画像技术将告警降噪率提升至80\%以上；
    \item 实际应用中助力用户实现超过85\%的告警降噪和分钟级威胁响应；
    \item 平均响应时间压缩至分钟级。
\end{itemize}

\subsection{绿盟科技：AI安全运营平台}

绿盟科技通过“风云卫AI安全能力平台”实现告警降噪，依托多个专业化智能体的分工协作，构建覆盖安全运营全流程的智能协同体系。

\textbf{效果数据}：
\begin{itemize}
    \item AI降噪率平均达到95\%以上；
    \item AI综合辅助研判准确率超过90\%；
    \item 依托自主响应可实现超过40\%的安全事件端到端自动化响应处置；
    \item 将应急响应从小时级压缩至分钟级。
\end{itemize}

\subsection{360数字安全：安全大模型告警降噪}

360通过“安全大模型+安全大脑”的结合实现告警降噪。360安全大模型由攻击检测、运营处置、追踪溯源、知识管理、数据保护、代码安全等六大专家子模型组成。

\textbf{效果数据}：
\begin{itemize}
    \item 某国有大型央企客户：告警降噪和事件自动化处置率提升90\%，事件研判与溯源调查时间大幅缩短96\%，人均处理效能提升300\%；
    \item 告警降噪效果提升近10倍，运营人效增长300\%；
    \item 某国有大型企业用户实现告警降噪效率提升900\%，事件研判时间缩短96\%。
\end{itemize}

\textbf{产品形态}：360本地安全大脑提供自动化告警降噪收敛、智能化关键信息聚合、交互式图谱拓线的威胁分析功能体系。

\subsection{亚信安全：AI XDR三重降噪}

亚信安全通过AI XDR平台实现告警降噪，采用基于机器学习的AI大数据三重降噪技术和智能研判功能。

\textbf{技术架构}：
\begin{itemize}
    \item \textbf{多维过滤机制}：通过高效的归并算法，将成千上万条重复、冗余的误报进行压缩；
    \item \textbf{客户本地数据自学习}：利用客户本地数据进行模型训练和自我学习，构建客户特定环境自适应威胁特性分析的智能模型；
    \item \textbf{AI大模型}：提供告警解读、告警降噪、基于安全处置生成指令等能力。
\end{itemize}

\textbf{效果数据}：
\begin{itemize}
    \item 98\%的噪声降减；
    \item 降低15\%-30\%安全运营总成本，增强30\%-50\%的高风险事件发现与闭环能力，事件响应效率提升50\%-70\%。
\end{itemize}

\section{横向对比分析}

\subsection{降噪效果对比}

\begin{center}
\begin{tabular}{>{\raggedright\arraybackslash}p{3cm}>{\raggedright\arraybackslash}p{5cm}>{\raggedright\arraybackslash}p{5cm}}
\toprule
\textbf{厂商} & \textbf{宣称降噪率} & \textbf{关键效果指标} \\
\midrule
奇安信 & 90\%-98\% & 1分钟事件定性，5分钟影响面评估 \\
深信服 & 95\%-99\% & 日均告警千万级→10+条 \\
腾讯云 & 97.7\% & 17万条→42条需关注 \\
微步在线 & 80\%-85\% & 分钟级响应 \\
绿盟科技 & 95\%以上 & 40\%+事件自动响应 \\
360 & 90\%-900\%效率提升 & 研判时间缩短96\% \\
亚信安全 & 98\% & 总成本降低15\%-30\% \\
阿里云 & 60\%-80\% & 分级管理+合并降噪 \\
\bottomrule
\end{tabular}
\end{center}

\subsection{技术路线分类}

根据技术架构的差异，可将上述厂商方案分为三类：

\textbf{第一类：平台+大模型深度融合型}（奇安信、深信服）
\begin{itemize}
    \item 以安全运营平台（NGSOC/XDR）为基座，深度融合自研安全大模型；
    \item 覆盖从告警接入、降噪、研判到处置的完整闭环；
    \item 强调“AI原生”而非“AI辅助”。
\end{itemize}

\textbf{第二类：Multi-Agent智能体型}（腾讯云、绿盟科技）
\begin{itemize}
    \item 采用多智能体分工协作架构；
    \item 将复杂分析拆解为多个子任务，由不同Agent分别完成；
    \item 强调“像安全专家一样主动调查”。
\end{itemize}

\textbf{第三类：情报+规则驱动型}（微步在线、阿里云）
\begin{itemize}
    \item 以威胁情报或规则引擎为核心驱动力；
    \item 通过IP画像、事件分组、合并降噪等实现告警削减；
    \item 更轻量，适合特定场景。
\end{itemize}

\subsection{国家级测试表现}

2024年，国家互联网应急中心（CNCERT）组织“人工智能技术赋能网络安全应用测试”，在“网络安全告警日志降噪”场景中，京东科技、绿盟科技、深信服、奇安信、安恒信息的团队取得较好成绩。2025年测试中，深信服、远江盛邦、奇安信的团队在该场景取得较好成绩。

这些测试结果表明，告警降噪已成为AI赋能安全运营的标杆场景，头部厂商在该领域已形成较强的技术积累。

\section{与“够用就行”理论体系的对话接口}

本报告梳理的工程实践与“基于生成论的网络安全分析能力调度系统”之间存在多个潜在的对话接口。

\subsection{降噪率的经济学追问}

当前所有厂商都在竞赛降噪率指标（90\%、95\%、97.7\%、99\%……）。然而，“够用就行”理论会追问一个根本问题：\textbf{“多高的降噪率才是‘够用’的？”}

当降噪率从95\%提升到99\%时，所需的额外研发投入、算力成本和模型复杂度增长了多少？这4\%的提升所能避免的安全损失又有多大？如果边际成本超过边际收益，那么继续追求更高的降噪率在经济学上是不合算的。“边际收益等于边际成本”原则为这场无止境的数字竞赛提供了一个理性的终点。

\subsection{优先级排序的统一框架}

厂商们通过复杂的模型（分诊、Agent、情报）来给告警打分排序，这本质上是在追求更准的“分类器”。而“够用就行”理论提供了一个更底层的视角：\textbf{所有这些努力，最终都可以统一到“成本-收益”的框架下}。

无论是AI计算的“危险分数”、情报给出的“威胁等级”，还是专家经验赋予的“优先级”，最终都可以也应该被转化为一个经济指标——分析这条告警的成本（$\kappa$）与漏掉它的预期损失（$\alpha A$）之间的比较。当$\alpha A > \kappa$时告警值得分析，反之则不值得。这一判据为各种复杂的优先级算法提供了一个统一的最优性标准。

\subsection{智能分诊的第一性原理}

奇安信的“智能分诊”、腾讯云的“告警触发与归并”、深信服的“威胁定性”等，都是工程上的优秀实践。“够用就行”理论为这些操作提供了第一性原理的解释：
\begin{itemize}
    \item \textbf{为什么能“自动过滤误报”？} 因为那些告警的A值（漏报损失）经评估后远低于分析成本（$\kappa$）；
    \item \textbf{为什么能“自动合并重复告警”？} 因为同源或同类的重复告警，其边际安全收益在递减，符合“同源关联优先”的数学原理；
    \item \textbf{为什么需要“智能分诊”？} 因为系统的分析能力（$C$）是有限的，必须通过分诊将资源分配给边际收益最高的告警，这正是“容量约束”和“贪心择优”的体现。
\end{itemize}

\subsection{理论指导工程优化的潜在方向}

“够用就行”理论可以为现有工程方案提供以下优化方向：
\begin{enumerate}
    \item \textbf{从“追求更高降噪率”转向“追求最优降噪率”}：不再以降噪率为唯一目标，而是以总成本（漏报损失+分析成本）最小化为目标；
    \item \textbf{引入经济参数作为模型输入}：将A值（资产损失估值）、$\kappa$（分析成本）等经济参数作为AI模型的输入特征，使优先级排序直接服务于成本-收益优化；
    \item \textbf{为降噪效果提供经济学评估}：不仅报告“降噪率X\%”，还报告“节省了Y元分析成本，带来了Z元漏报风险”，使安全决策可量化、可审计。
\end{enumerate}

\section{结论}

本报告系统梳理了国内八家主流安全厂商的告警降噪技术方案。总体而言，当前业界方案呈现出以下特征：

\begin{enumerate}
    \item \textbf{技术范式高度趋同}：以“AI大模型+安全平台”为核心架构，强调自动化、智能化的告警处理；
    \item \textbf{效果数据普遍亮眼}：多数厂商宣称降噪率达到90\%以上，部分案例甚至达到99\%；
    \item \textbf{缺乏统一的最优性标准}：各厂商以“降噪率”为竞争指标，但无人追问“多高的降噪率才是够用的”；
    \item \textbf{工程先行，理论滞后}：业界在工程实践上走在前面，但缺乏统一的理论框架来解释“为什么有效”以及“何时该停止优化”。
\end{enumerate}

“基于生成论的网络安全分析能力调度系统”恰好填补了这一理论空白。它提供了一个统一的经济学框架，将告警分级、资源分配、优先级排序等分散的工程问题统一到“边际收益等于边际成本”的数学判据之下。这一理论不仅可以解释现有方案为什么有效，更能指引未来优化的方向——不是无休止地追求更复杂的模型和更高的降噪率，而是去无限逼近那个由经济学原理所定义的理论最优平衡点。

\vspace{5mm}
\noindent\textbf{报告说明}：
\begin{itemize}
    \item 本报告所有数据均来源于各厂商官方网站、公开新闻稿、技术白皮书及国家级测试报告；
    \item 各厂商宣称的降噪率因测试环境、数据来源和计算口径不同，不宜直接横向比较；
    \item 本报告旨在提供工程实践的参考框架，不构成对任何厂商产品的推荐或评价。
\end{itemize}

\end{document}